\documentclass{article}
\usepackage[T2A]{fontenc}
\usepackage[utf8]{inputenc} % utf8
\usepackage{amsthm}
\usepackage{amsmath}
\usepackage{amssymb}
\usepackage{amsfonts}
\usepackage{mathrsfs}
\usepackage[12pt]{extsizes}
\usepackage{fancyvrb}
\usepackage{indentfirst}
\usepackage[
  left=2cm, right=2cm, top=2cm, bottom=2cm, headsep=0.2cm, footskip=0.6cm, bindingoffset=0cm
]{geometry}
\usepackage[english,russian]{babel}

\begin{document}
\section*{Вариант 6}

Общая оценка последовательности $S$ строится как сумма оценок каждого входящего в неё вектора $X_{i}$
\begin{equation}
\begin{split}
O(A_{0},F_{j},S) &= \sum\limits_{i=1}^{\text{длина}(S)}O(A_{0},F_{j},X_{i}) = \\
&= \sum\limits_{i=1}^{\text{длина}(S)}(c_{1}\cdot N_{1}(X_{i},A_{0},F_{j})+c_{2}\cdot N_{2}(X_{i},A_{0},F_{j})+c_{3}\cdot N_{3}(X_{i},A_{0},F_{j}))
\end{split}
\end{equation}
где: $X_{i}$ "--- $i$-й набор последовательности $S$; $c_{1}-c_{3}$ "--- нормализующие константы; $N_{1}-N_{3}$ "--- параметры различающей активности.

Функция различия для (1) переопределяется следующим образом:
\begin{equation}
R(g, X_{i}, A_{0}, F_{j}) =
\begin{cases}
0, & \text{если выходы элементов } g \text{ множества } U, \\
   & \text{которые принадлежат классу } F_{j}, \\
   & \text{одинаковы после подачи набора } X_{i}; \\
1, & \text{иначе.}
\end{cases}
\end{equation}

\end{document}
